\documentclass[11pt]{article}
\usepackage[T1]{fontenc}
\usepackage{lmodern}
\usepackage[margin=1in]{geometry}
\usepackage{amsmath,amssymb,amsthm,mathtools}
\usepackage{booktabs,microtype}
\usepackage{xurl}
\usepackage[hidelinks]{hyperref}
\newtheorem{theorem}{Theorem}
\newtheorem{lemma}[theorem]{Lemma}
\newtheorem{proposition}[theorem]{Proposition}
\theoremstyle{definition}
\newtheorem{definition}[theorem]{Definition}
\theoremstyle{remark}
\newtheorem{remark}[theorem]{Remark}
\newcommand{\R}{\mathbb R}
\newcommand{\T}{\mathbb T}
\newcommand{\one}{\mathbf 1}
\newcommand{\Pp}{\mathbb P}
\newcommand{\norm}[1]{\left\lVert#1\right\rVert}
\newcommand{\ip}[2]{\left\langle#1,#2\right\rangle}
\newcommand{\asin}{\operatorname{arcsin}}
\title{Random cubic graphs admit nonsynchronized stable Kuramoto equilibria}
\author{}
\date{}
\begin{document}
\maketitle
\vspace{-2.5em}
\begin{abstract}
For a uniformly random labelled simple cubic graph on an even number of
vertices, we give a proof claim that the probability that every local minimum
of the homogeneous attractive Kuramoto energy is synchronized tends to zero.
This is the opposite of the limit proposed in problem MD-06 of
\emph{OpenProblemsInNLA} and Conjecture~3 of
\emph{Randomstrasse101: Open Problems of 2024}.
A two-by-two sign pattern on a cycle of length divisible by four, together
with a decaying nonlinear profile on its outward binary trees, has uniformly
positive edge cosines and an exponentially small boundary gradient.
A Laplacian spectral gap converts this truncated profile into an exact local
minimum. Fixed-length cycle counts and the spectral gap of random regular
graphs then establish the claimed probability limit. The analytic argument
does not depend on numerical optimization. Supplementary rational existence certificates and a symbolic finite
construction are checked with exact arithmetic.
\end{abstract}

\noindent\textbf{Review status.} Complete negative-resolution proof claim;
independent mathematical review is pending. No external peer review or formal
proof-assistant verification is asserted.

\section{Statement and scope}

Let $G=(V,E)$ be a finite simple undirected graph, and let
$\T=\R/(2\pi\mathbb Z)$. Define
\begin{equation}\label{eq:energy}
 E_G(\theta)=\sum_{\{u,v\}\in E}\bigl(1-\cos(\theta_u-\theta_v)\bigr),
 \qquad \theta\in\T^{V}.
\end{equation}
A phase vector is \emph{synchronized} when all coordinates agree modulo
$2\pi$. Local minima use the usual topology on the phase torus. Because a
common phase rotation is a symmetry, a local minimum need not be strict on
the entire torus.

Write $L_G$ for the unweighted graph Laplacian, and $\lambda_2(L_G)$ for its
second-smallest eigenvalue. For connected $G$, this eigenvalue is positive.
The subspace $\one^\perp$ is the orthogonal complement of the constant vector
in $\R^{V}$.

\begin{theorem}\label{thm:main}
Let $G_n$ be uniform over the labelled simple $3$-regular graphs on
$\{1,\ldots,n\}$, for even $n\ge4$. Then
\begin{equation}\label{eq:main}
 \lim_{\substack{n\to\infty\\ n\ \mathrm{even}}}
 \Pp\bigl(\text{every local minimum of }E_{G_n}
                \text{ is synchronized}\bigr)=0.
\end{equation}
More strongly, with probability tending to one, there is a nonsynchronized
critical point with all edge cosines bounded below by the same positive
absolute constant. Its Hessian is positive definite on $\one^\perp$.
\end{theorem}

The graph model, energy, phase domain, and quantification over \emph{every}
local minimum match the canonical MD-06 entry~\cite{repo}. Its source is
Definition~2.1 and Conjecture~3 of Bandeira, Kireeva, Maillard, and
R\"odder~\cite{randomstrasse}. Therefore, if independently verified,
Theorem~\ref{thm:main} gives a full negative resolution, not just a family
of exceptional finite graphs.

Stable nonsynchronized patterns on finite cubic graphs are not new.
DeVille and Ermentrout~\cite{deville} studied them analytically for specific
graphs and numerically across finite graph samples. Their Section~IV
explicitly identifies a vanishing probability of having no stable pattern
as a trend suggested by their data. The argument below is directed at the
asymptotic assertion for the uniform labelled simple model; no priority
claim about finite examples or the numerical phenomenon is made.

\section{An exact minimum from a small gradient}

We work with real lifts of the angles. The energy and its derivatives are
periodic in each coordinate. Direct differentiation gives
\begin{align}
 (\nabla E_G(\phi))_v
   &=\sum_{w\sim v}\sin(\phi_v-\phi_w),\label{eq:gradient}\\
 z^{\mathsf T}\nabla^2E_G(\phi)z
   &=\sum_{\{u,v\}\in E}\cos(\phi_u-\phi_v)(z_u-z_v)^2.
   \label{eq:hessian}
\end{align}
The gradient has zero coordinate sum, and the Hessian annihilates $\one$.

\begin{lemma}[Small-gradient certificate]\label{lem:certificate}
Suppose $G$ is connected, $\lambda_2(L_G)\ge\gamma>0$, and
$\phi\in\R^V$ satisfies
\begin{equation}\label{eq:certificate-hyp}
 \cos(\phi_u-\phi_v)\ge c>0\quad(\{u,v\}\in E),
 \qquad
 \norm{\nabla E_G(\phi)}_2<\frac{c^2\gamma}{4\sqrt2}.
\end{equation}
There is $h_*\in\one^\perp$, with $\norm{h_*}_2<c/(2\sqrt2)$, such that
$\theta_* =\phi+h_*$ is an exact critical point. Furthermore,
\begin{align}
 \cos((\theta_*)_u-(\theta_*)_v)&\ge c/2
       &&(\{u,v\}\in E),\label{eq:positive-cos}\\
 z^{\mathsf T}\nabla^2E_G(\theta_*)z
       &\ge(c\gamma/2)\norm z_2^2
       &&(z\in\one^\perp).\label{eq:positive-hessian}
\end{align}
Thus $\theta_*$ is a strict local minimum modulo common rotation and a local
minimum on the phase torus.
\end{lemma}

\begin{proof}
Set $\rho=c/(2\sqrt2)$ and
$B=\{h\in\one^\perp:\norm h_2\le\rho\}$.
For each edge and each $h\in B$,
$|h_u-h_v|\le\sqrt2\norm h_2\le c/2$.
Since cosine is $1$-Lipschitz, all weights in \eqref{eq:hessian} at
$\phi+h$ are at least $c/2$. Throughout $B$ this implies
\begin{equation}\label{eq:strongconvex}
 z^{\mathsf T}\nabla^2E_G(\phi+h)z
  \ge(c/2)z^{\mathsf T}L_Gz
  \ge\mu\norm z_2^2\quad(z\in\one^\perp),\qquad \mu=c\gamma/2.
\end{equation}
The continuous function $h\mapsto E_G(\phi+h)$ attains a minimum on the
compact set $B$. At a boundary point, $\norm h_2=\rho$, and
\begin{align*}
 \ip h{\nabla E_G(\phi+h)}
 &=\ip h{\nabla E_G(\phi)}
   +\int_0^1h^{\mathsf T}\nabla^2E_G(\phi+th)h\,dt\\
 &\ge-\rho\norm{\nabla E_G(\phi)}_2+\mu\rho^2>0.
\end{align*}
A small move in the direction $-h$ decreases the energy, so a minimizer
$h_*$ is in the relative interior of $B$. At this point the gradient is
orthogonal to $\one^\perp$. It is also orthogonal to $\one$, hence zero.
The weight bound and \eqref{eq:strongconvex} prove
\eqref{eq:positive-cos}--\eqref{eq:positive-hessian} and strict local
minimality on the slice. Any nearby perturbation splits into a mean-zero
perturbation and a common rotation, proving local minimality on the torus.
\end{proof}

\section{A cycle with outward binary trees}

\subsection{A scalar phase profile}

For each integer $R\ge4$, define
\begin{equation}\label{eq:F}
 F_R(t)=3\asin t+2\sum_{j=1}^{R-1}\asin(t/2^j),\qquad 0\le t\le1.
\end{equation}

\begin{lemma}\label{lem:profile}
There is a unique $t_R\in(0,1)$ satisfying $F_R(t_R)=2\pi$.
The roots decrease with $R$, and $t_R>1/2$. Hence
\begin{equation}\label{eq:c0}
 c_0:=\sqrt{1-t_4^2}>0,\qquad \sqrt{1-t_R^2}\ge c_0\quad(R\ge4).
\end{equation}
Set
\begin{equation}\label{eq:profile}
 \delta_j=\asin(t_R/2^j),\qquad
 a_j=\sum_{s=j}^{R-1}\delta_s\ (0\le j<R),\qquad a_R=0.
\end{equation}
Then
\begin{equation}\label{eq:rootbalance}
 2a_0+\delta_0=2\pi,\qquad
 \sin(2a_0)=-t_R,\qquad
 \cos(2a_0)=\sqrt{1-t_R^2}.
\end{equation}
\end{lemma}

\begin{proof}
Each $F_R$ is continuous and strictly increasing, with $F_R(0)=0$.
Moreover,
\[
 F_4(1)=\frac{11\pi}{6}+2\asin(1/4)+2\asin(1/8)
       >\frac{11\pi}{6}+\frac34>2\pi,
\]
where the last inequality follows from $\pi<4$.
Since $F_R\ge F_4$, existence and uniqueness follow by the intermediate
value theorem. Increasing $R$ increases $F_R(t)$ for $t>0$, so the roots
decrease.

Convexity of arcsine on $[0,1]$ gives $\asin x\le\pi x/2$.
Using $3\asin(1/2)=\pi/2$, we obtain
\[
 F_R(1/2)\le\frac\pi2+\pi\sum_{j=1}^{\infty}2^{-j-1}=\pi<2\pi.
\]
This proves $t_R>1/2$. Finally, the identity
$F_R(t_R)=2a_0+\delta_0$ gives \eqref{eq:rootbalance}.
\end{proof}

\begin{remark}[Explicit constants]\label{rem:constants}
The constant in \eqref{eq:c0} satisfies $c_0>1/16$. To see this, put
$t_*=\sqrt{255}/16$. The elementary inequalities
$\asin x\ge x+x^3/6$ and
$\asin(1/16)<1/\sqrt{255}<1/15$ give
\[
 F_4(t_*)>
 \frac{3\pi}{2}-\frac15+\frac{7t_*}{4}+\frac{t_*^3}{24}
 >\frac{3\pi}{2}+\frac{151}{96}>2\pi.
\]
Here $t_*>99/100$, $t_*^3>97/100$, and
$151/96>11/7>\pi/2$ were used. Thus $t_4<t_*$, proving the claim.
Consequently, the final edge-cosine bound can be taken as $1/32$ and,
when $\gamma=1/10$, the Hessian bound on $\one^\perp$ as $(1/320)I$.
In that case the explicit choice
$R=29+\lceil\log_2\ell\rceil$ satisfies \eqref{eq:chooseR}, since
$32\ell/((1/16)^4(1/10)^2)=209\,715\,200\ell<2^{28}\ell$.
These conservative constants are not needed for the probability limit.
\end{remark}

\begin{definition}\label{def:clean}
A cycle $C$ in a cubic graph has a \emph{clean radius-$R$ neighborhood} when
the subgraph induced by the vertices at distance at most $R$ from $C$ is
unicyclic, with $C$ as its unique cycle.
\end{definition}

For a clean neighborhood of an $\ell$-cycle, every cycle vertex has one
outward child. Each vertex at depth $1,\ldots,R-1$ has two children. The
attached rooted trees are disjoint, and their depth-$j$ level has
$\ell2^{j-1}$ vertices, for $1\le j\le R$.

\begin{proposition}[Localized approximate equilibrium]\label{prop:localized}
Let $C=(v_0,\ldots,v_{\ell-1})$ have length divisible by four and a clean
radius-$R$ neighborhood in a cubic graph, where $R\ge4$. There is
$\phi\in\R^V$ such that
\begin{equation}\label{eq:localized-estimates}
 \min_{\{u,v\}\in E}\cos(\phi_u-\phi_v)\ge c_0,
 \qquad \norm{\nabla E_G(\phi)}_2^2=\ell t_R^2 2^{1-R}.
\end{equation}
On an opposite-sign cycle edge in the construction, the distance of its
phase difference from $2\pi\mathbb Z$ is $\delta_0>\pi/6$.
\end{proposition}

\begin{proof}
Give the cycle roots signs
\[
 (\varepsilon_0,\ldots,\varepsilon_{\ell-1})
   =(+1,+1,-1,-1,+1,+1,-1,-1,\ldots).
\]
Every root has one same-sign and one opposite-sign cycle neighbor.
Assign phase $\varepsilon_i a_j$ to each depth-$j$ vertex in the tree rooted
at $v_i$, for $j<R$, including $j=0$ for the root. Give every other vertex
phase zero, including all depth-$R$ vertices.

At a root of sign $\varepsilon$, the three terms in the gradient are
\[
 0,\qquad \sin(2\varepsilon a_0)=-\varepsilon t_R,
 \qquad \sin\bigl(\varepsilon(a_0-a_1)\bigr)=\varepsilon t_R.
\]
They cancel. At a depth-$j$ vertex, $1\le j<R$, the parent contributes
$-\varepsilon t_R/2^{j-1}$ and each of its two children contributes
$\varepsilon t_R/2^j$. These also cancel. A depth-$R$ vertex has phase zero,
one nonzero-phase parent, and two zero-phase neighbors. Its gradient is
$-\varepsilon t_R/2^{R-1}$. Beyond depth $R$ every gradient coordinate is
zero. The number of boundary vertices is $\ell2^{R-1}$, which proves the
norm identity.

Opposite-sign cycle edges have cosine
$\cos(2a_0)=\sqrt{1-t_R^2}\ge c_0$; same-sign cycle edges have cosine one.
A tree edge from depth $j$ to $j+1$ has cosine
$\cos\delta_j\ge\cos\delta_0\ge c_0$. All other edges have phase difference
zero. Lastly, a cross-sign cycle edge has difference
$\pm(2\pi-\delta_0)$, where $\delta_0=\asin t_R>\pi/6$.
\end{proof}

\begin{theorem}[Deterministic obstruction]\label{thm:deterministic}
Fix $\gamma>0$ and $\ell\in4\mathbb N$, and choose an integer $R\ge4$ with
\begin{equation}\label{eq:chooseR}
 2^{R-1}>\frac{32\ell}{c_0^4\gamma^2}.
\end{equation}
Any connected cubic graph $G$ with $\lambda_2(L_G)\ge\gamma$ and a clean
radius-$R$ neighborhood of an $\ell$-cycle has a nonsynchronized local
minimum. At this minimum every edge cosine is at least $c_0/2$ and the
Hessian on $\one^\perp$ is bounded below by $(c_0\gamma/2)I$.
\end{theorem}

\begin{proof}
Proposition~\ref{prop:localized} and $t_R<1$ give
\[
 \norm{\nabla E_G(\phi)}_2=t_R\sqrt{\ell2^{1-R}}
            <\frac{c_0^2\gamma}{4\sqrt2}.
\]
Apply Lemma~\ref{lem:certificate} with $c=c_0$. Every edge difference changes
by less than $c_0/2\le1/2$. On a cross-sign cycle edge, the original distance
from $2\pi\mathbb Z$ is $\delta_0>\pi/6>1/2$. Distance to a fixed closed
set is $1$-Lipschitz, so the new difference is not a multiple of $2\pi$.
The exact local minimum is therefore nonsynchronized. Its other properties
are the conclusions of the certificate lemma.
\end{proof}

\section{Uniform random simple cubic graphs}

We use two established results in their uniform labelled simple-graph forms.
Bordenave's version of Friedman's theorem
\cite[Theorem~1, printed p.~2]{bordenave} bounds the nontrivial adjacency
eigenvalues of a random $d$-regular graph by
$2\sqrt{d-1}+o(1)$ in absolute value, with probability tending to one.
For $d=3$, since $3-2\sqrt2>1/10$, it follows that
\begin{equation}\label{eq:gap-random}
 \Pp\bigl(\lambda_2(L_{G_n})\ge1/10\bigr)\longrightarrow1.
\end{equation}
The event in \eqref{eq:gap-random} also implies connectivity.

Let $C_\ell(G_n)$ count simple cycles of length $\ell$. Johnson's
Poisson approximation \cite[Theorem~11, printed p.~12]{johnson} gives,
for each fixed $\ell\ge3$,
\begin{equation}\label{eq:poisson}
 C_\ell(G_n)\xrightarrow{\ d\ }\operatorname{Poisson}(\mu_\ell),
 \qquad \mu_\ell=\frac{2^\ell}{2\ell}.
\end{equation}
The simple-graph convention is stated in Section~2, printed p.~3, of that
paper. In particular, $\Pp(C_\ell(G_n)=0)\to e^{-\mu_\ell}$.

\begin{lemma}[Clean fixed neighborhoods]\label{lem:clean-random}
For fixed $\ell\ge3$ and $R\ge1$, the probability that $G_n$ contains an
$\ell$-cycle without a clean radius-$R$ neighborhood is
$O_{\ell,R}(n^{-1})$.
\end{lemma}

\begin{proof}
A radius-$R$ neighborhood of an $\ell$-cycle in a cubic graph has at most
$M=\ell2^R$ vertices. If it is not unicyclic, it is connected and has at
least two independent cycles. Repeatedly remove degree-one vertices. The
remaining graph $H$ is nonempty, has minimum degree at least two, maximum
degree at most three, at most $M$ vertices, and
$e(H)-v(H)\ge1$. Removing a leaf preserves $e-v$.

There are finitely many such isomorphism types, with a number depending
only on $\ell,R$. For any fixed labelled embedding of one of these types,
Johnson's Proposition~1(a), printed p.~3, gives
\[
 \Pp(H\subseteq G_n)\le c_1\,2^{e(H)} n^{-e(H)}
\]
for sufficiently large $n$. Its degree and size conditions hold because
$d=3$ and $H$ is fixed. There are at most $n^{v(H)}$ embeddings of a type,
so the union bound over all types is $O_{\ell,R}(n^{-1})$.
\end{proof}

\begin{proof}[Proof of Theorem~\ref{thm:main}]
Write $S_n$ for the event that every local minimum is synchronized.
Fix any $\ell\in4\mathbb N$, set $\gamma=1/10$, and choose a fixed $R$
satisfying \eqref{eq:chooseR}. The deterministic obstruction implies
\begin{align*}
 \Pp(S_n)
 &\le \Pp\bigl(\lambda_2(L_{G_n})<1/10\bigr)
  +O_{\ell,R}(n^{-1})+\Pp\bigl(C_\ell(G_n)=0\bigr).
\end{align*}
Indeed, if the gap condition holds, no $\ell$-cycle has a dirty neighborhood,
and an $\ell$-cycle exists, then a nonsynchronized local minimum exists.
Equations~\eqref{eq:gap-random}--\eqref{eq:poisson} yield
\begin{equation}\label{eq:limsup}
 \limsup_{\substack{n\to\infty\\n\ \mathrm{even}}}\Pp(S_n)
       \le\exp\!\left(-\frac{2^\ell}{2\ell}\right).
\end{equation}
This is true for \emph{every fixed} multiple of four. Let $\ell$ tend to
infinity after taking the limit superior; the right side tends to zero.
This proves \eqref{eq:main}.

The same bound applies to the event of having no nonsynchronized minimum
with edge cosines at least $c_0/2$ and Hessian lower bound $(c_0/20)I$ on
$\one^\perp$. The constants are independent of $\ell$, proving the stronger
assertion.
\end{proof}

\begin{remark}[Quantifiers and scope]\label{rem:scope}
The events in the union bound need not be independent. In each use of a
random-graph limit, both $\ell$ and $R$ are fixed before $n$ tends to
infinity. No estimates uniform in growing cycle length or neighborhood
radius are required. The radius bound is conservative, so this proof does
not supply a useful numerical threshold in $n$. Nor does it assert that a
uniform random initial phase converges to a nonsynchronized minimum with
probability tending to one: existence of a local minimum is the target.
\end{remark}

\section{Supplementary exact finite certificates}

\subsection{Rational approximate phases}

Two supplied rational approximate phase vectors, on connected simple cubic
graphs with $20$ and $500$ vertices, satisfy the hypotheses of
Lemma~\ref{lem:certificate}. The standalone verifier uses only integer and
rational arithmetic for acceptance. It proves existence of nearby exact
minima, not that the supplied rational vectors themselves are stationary.
These finite checks are supplementary, not premises of the asymptotic proof.

For a connected graph with $n$ vertices and diameter $D$,
\begin{equation}\label{eq:diamgap}
 \lambda_2(L_G)\ge\frac{2}{(n-1)D}>\frac1{nD}.
\end{equation}
To see this, for $z\in\one^\perp$ use a shortest path for each vertex pair:
\[
 n\norm z_2^2=\sum_{u<v}(z_u-z_v)^2
 \le\frac{n(n-1)}2 D\sum_{\{u,v\}\in E}(z_u-z_v)^2.
\]
The verifier computes $D$ by exact breadth-first searches and uses
$\gamma=1/(nD)$.

For a rational edge difference $x$, it evaluates
\[
 S_N(x)=\sum_{j=0}^N\frac{(-1)^jx^{2j+1}}{(2j+1)!},\qquad
 C_N(x)=\sum_{j=0}^N\frac{(-1)^jx^{2j}}{(2j)!}.
\]
Including the next zero Taylor coefficient, Lagrange's remainder theorem
gives
\[
 |\sin x-S_N(x)|\le\frac{|x|^{2N+3}}{(2N+3)!},\qquad
 |\cos x-C_N(x)|\le\frac{|x|^{2N+2}}{(2N+2)!}.
\]
These rational enclosures hold for every real $x$; they do not require an
alternating-series monotonicity assumption. The supplied runs use $N=40$.
Summing sine intervals at each vertex gives an upper bound $g_2$ on the
squared gradient norm. The decisive test is the exact rational inequality
\begin{equation}\label{eq:squaredtest}
 g_2<\frac{c^4\gamma^2}{32}.
\end{equation}
A separate test finds an edge with original absolute real phase difference
between $1$ and $5$. Its change is at most $c/2\le1/2$, so it stays strictly
between $0$ and $2\pi$ in absolute value, since $2\pi>6$.

\begin{center}
\small
\begin{tabular}{@{}rrrrl@{}}
\toprule
$n$ & $D$ & $c$ & $\gamma$ & verified conservative bound\\
\midrule
20 & 5 & $1/4$ & $1/100$ & $g_2<10^{-22}$\\
500 & 12 & $1/8$ & $1/6000$ & $g_2<10^{-21}$\\
\bottomrule
\end{tabular}
\end{center}
The right sides of \eqref{eq:squaredtest} are, respectively,
$1/81\,920\,000$ and $1/4\,718\,592\,000\,000$.
All acceptance comparisons use \texttt{fractions.Fraction}; decimal log
entries are descriptive only. Floating-point optimization and Newton
refinement were used only to propose rational inputs to the verifier.

\subsection{A symbolic finite realization}

There is also a finite realization of the tree profile with an exactly zero
gradient, requiring no perturbation argument. Start with the four cycle roots
of signs $+,+,-,-$, and build their outward trees through depth $R-1$.
At depth $R$, pair each prospective leaf of the first positive tree with a
leaf of the first negative tree, and do the same for the other two trees.
Identify each opposite-sign pair into a single vertex of phase zero.
Every such vertex has two parents of opposite phases and therefore has zero
net sine force. Match the resulting zero-phase vertices between the two
pairing groups, giving each one a third edge.

The resulting connected simple cubic graph has $3\cdot2^R$ vertices.
The phase expressions are integer linear combinations of
$\delta_0,\ldots,\delta_{R-1}$. Every edge difference is one of
$0$, $\pm\delta_j$, or $\pm2a_0=\pm(2\pi-\delta_0)$.
Consequently, each sine contribution is an exact rational multiple of $t_R$,
and all gradient coordinates vanish. Every edge cosine is positive, so
\eqref{eq:hessian} proves strict local minimality modulo rotation. The
cross-root gap is $\delta_0>\pi/6$, so the state is nonsynchronized.

The supplied \texttt{analytic\_gadget.py} constructs this graph and checks
its degree, simplicity, connectivity, symbolic edge types, and all sine-flow
balances using integers and rational numbers. At $R=4$ it has $48$ vertices
and $72$ edges; the accompanying tests also check $R=5$. This example
illustrates the same nonlinear profile but, like the other finite examples,
is not a substitute for the random-graph argument.

\section{Points for independent review}

The proof has four distinct steps: a scalar equation supplies the positive
cosine margin; the sign pattern and tree profile cancel all interior
forces; compact-ball minimization corrects the boundary residual despite
arbitrary exterior connections; and the fixed-cycle Poisson limit yields
\eqref{eq:limsup}. In particular, numerical stationarity tolerances, finite
examples, and unproved independence assumptions are not used to infer an
asymptotic result. Independent review should check the analytic chain and
the exact hypotheses of the two random-graph inputs, rather than only rerun
the finite verifier.

\begin{thebibliography}{9}
\small
\setlength{\itemsep}{4pt}
\bibitem{repo}
\emph{OpenProblemsInNLA}, canonical entry MD-06,
``Global synchronization of a random cubic graph.''
Entry checked September 11, 2026; its recorded source-check date is
September 10, 2026.
\href{https://github.com/ajt60gaibb/OpenProblemsInNLA/tree/main/matrix-discrepancy-and-optimization/MD-06}{Repository entry MD-06}.

\bibitem{randomstrasse}
A.~S. Bandeira, A. Kireeva, A. Maillard, and A. R\"odder,
\emph{Randomstrasse101: Open Problems of 2024},
arXiv:2504.20539v1 (2025), Definition~2.1 and Conjecture~3, printed p.~3.
\url{https://arxiv.org/abs/2504.20539v1}.

\bibitem{deville}
L. DeVille and B. Ermentrout,
\emph{Phase-locked patterns of the Kuramoto model on 3-regular graphs},
Chaos \textbf{26} (2016), 094820.
DOI: \href{https://doi.org/10.1063/1.4961064}{10.1063/1.4961064}.
Preprint: \url{https://arxiv.org/abs/1512.06140}.
Relevant locators: Sections~I.E, II.A, III.G, and IV of the arXiv HTML version.

\bibitem{bordenave}
C. Bordenave,
\emph{A new proof of Friedman's second eigenvalue theorem and its extension
to random lifts}, arXiv:1502.04482.
Theorem~1 and the labelled simple-graph model, printed p.~2.
\url{https://arxiv.org/abs/1502.04482}.

\bibitem{johnson}
T. Johnson,
\emph{Exchangeable pairs, switchings, and random regular graphs},
arXiv:1112.0704v5.
Section~2 and Proposition~1(a), printed p.~3;
Theorem~11, printed p.~12.
\url{https://arxiv.org/abs/1112.0704v5}.
\end{thebibliography}
\end{document}
